\section*{Выводы}

В ходе выполнения работы было проведено комплексное моделирование малошумящего транзисторного СВЧ-усилителя на базе биполярного транзистора NE68133. 

\begin{itemize}
    \item На этапе линейного анализа была успешно решена проблема потенциальной неустойчивости транзистора на заданной рабочей частоте 2.5 ГГц. Путем введения стабилизирующей индуктивности в цепь эмиттера была обеспечена абсолютная устойчивость каскада. С использованием диаграммы Смита были синтезированы входная и выходная согласующие цепи, позволившие достичь требуемого баланса между коэффициентом усиления и минимальным коэффициентом шума.
    
    \item Проведен анализ по постоянному току (DC анализ) нелинейной SPICE-модели транзистора. Получено семейство выходных вольтамперных характеристик, необходимых для корректного выбора рабочей точки и оценки выходной мощности схемы.
    
    \item Расчет каскада методами AC и гармонического баланса (Harmonic Balance) позволил верифицировать малосигнальные параметры усилителя и исследовать его нелинейные свойства. Был определен динамический диапазон устройства: установлено, что точка децибельной компрессии (сжатие усиления на 1 дБ) достигается при уровне входного сигнала -8 дБм.
    
    \item Спроектирована геометрическая трехмерная модель выходного фильтра на базе круглого диэлектрического резонатора и микрополосковых линий. В связи с возникшими сложностями при конфигурации портов возбуждения, электромагнитный расчет фильтра и его итоговая интеграция в схему усилителя не были выполнены, что оставляет эту задачу для дальнейшей доработки.
\end{itemize}

В целом, базовые цели проектирования транзисторного каскада достигнуты: освоены методы согласования СВЧ-цепей с использованием диаграммы Смита, а также получены практические навыки работы с линейными и нелинейными моделями в современных САПР.